\subsubsection{Fidelity and Decoherence}\label{sec:fidelity-and-decoherence}

Imagine we start with a gate that takes time $t_{\text{gate}}$ to execute. During this time, decoherence attacks our qubit. Therefore, the longer the gate lasts, the more time decoherence has to damage the qubit's state. \hl{More decoherence means more errors, which means lower fidelity.}

\highspace
Precisely, the fidelity measures how accurately a physical gate reproduces the ideal gate (\autopageref{sec:fidelity-of-a-quantum-gate}). One important source of infidelity is decoherence (\autopageref{def:coherence}). Since a gate requires a \textbf{finite execution time} $t_{\text{gate}}$, the qubit is exposed to environmental noise while the gate is running. The longer the gate duration compared to the \textbf{decoherence time} $t_{\text{dec}}$, the higher the probability that decoherence introduces errors. Consequently, \textbf{longer gates tend to have lower fidelity}. A \textbf{simple approximation} is:
\begin{equation}
    p_{\text{gate}} \approx \frac{t_{\text{gate}}}{t_{\text{dec}}}
\end{equation}
Where:
\begin{itemize}
    \item $t_{\text{dec}}$ is the \definition{Decoherence Time}. It measures \hl{how long a qubit can preserve its quantum state before the environment destroys it.} For superconducting qubits, usually the decoherence time is on the order of microseconds ($\mu s$):
    \begin{equation*}
        t_{\text{dec}} \approx 200 \mu s
    \end{equation*}


    \item $t_{\text{gate}}$ is the \definition{Gate Time}. It measures \hl{how long a quantum gate takes to execute.} Generally, $X$ gates take around 20 ns, $H$ gates around 30 ns, and CNOT gates around 300 ns. The important assumption is $t_{\text{gate}} \ll t_{\text{dec}}$, because otherwise the qubit would lose coherence while the gate is still running, leading to very low fidelity.
    

    \item $p_{\text{gate}}$ is the \definition{Gate Error Probability}. It estimates \hl{the likelihood that a gate operation fails due to decoherence.} The longer the gate duration relative to the decoherence time, the higher the error probability, which reduces the overall fidelity of the quantum operation.
\end{itemize}
For example, suppose $t_{\text{dec}} = 100 \mu s$ and $t_{\text{gate}} = 1 \mu s$. Then:
\begin{equation*}
    p_{\text{gate}} \approx \frac{1 \mu s}{100 \mu s} = 0.01
\end{equation*}
This means that the gate has a 1\% chance of introducing an error due to decoherence, which corresponds to a fidelity of approximately 99\%. If we can reduce the gate time to 0.5 $\mu s$, then:
\begin{equation*}
    p_{\text{gate}} \approx \frac{0.5 \mu s}{100 \mu s} = 0.005
\end{equation*}
This means the gate error probability drops to 0.5\%, improving the fidelity to approximately 99.5\%. Therefore, \hl{optimizing gate times is crucial for enhancing the fidelity of quantum operations.}

\highspace
\textcolor{Green3}{\faIcon[regular]{star} \textbf{Quality Ratio.}} The \definition{Quality Ratio} is a metric that answers the question: \emph{\textbf{how many gate operations can we perform before decoherence becomes significant?}} It is defined as:
\begin{equation}
    \text{Quality Ratio} = \frac{t_{\text{dec}}}{t_{\text{gate}}}
\end{equation}
This is simply the inverse of the gate error probability:
\begin{equation*}
    \text{Quality Ratio} = \frac{1}{p_{\text{gate}}}
\end{equation*}
A \hl{higher quality ratio indicates that we can perform more gate operations before decoherence significantly degrades the qubit's state.} For instance, if $t_{\text{dec}} = 100 \mu s$ and $t_{\text{gate}} = 1 \mu s$, then the quality ratio is:
\begin{equation*}
    \text{Quality Ratio} = \frac{100 \mu s}{1 \mu s} = 100
\end{equation*}
This means we can perform approximately 100 gate operations before decoherence significantly impacts the qubit's state. If we reduce the gate time to 0.5 $\mu s$, the quality ratio increases to:
\begin{equation*}
    \text{Quality Ratio} = \frac{100 \mu s}{0.5 \mu s} = 200
\end{equation*}
This means we can perform approximately 200 gate operations before decoherence becomes significant, effectively doubling the number of operations we can execute with high fidelity. Therefore, optimizing gate times not only reduces error probabilities but also increases the quality ratio, allowing for more complex quantum computations before decoherence limits performance. In general, a quality ratio of $10^{3} - 10^{4}$ is the target for fault-tolerant quantum computing, as it allows for sufficient error correction to mitigate the effects of decoherence.

\highspace
\textcolor{Green3}{\faIcon{check-circle} \textbf{Updated Fidelity Formula.}} The \definition{Fidelity} of a quantum gate can be \textbf{approximated by considering the gate error probability due to decoherence}. If we assume that the only source of infidelity is decoherence, then the fidelity can be expressed as:
\begin{equation}
    F_{\text{gate}} \approx 1 - p_{\text{gate}} = 1 - \frac{t_{\text{gate}}}{t_{\text{dec}}}
\end{equation}
This formula indicates that the fidelity decreases linearly with the ratio of gate time to decoherence time. For example, if $t_{\text{dec}} = 100 \mu s$ and $t_{\text{gate}} = 1 \mu s$, then the fidelity is:
\begin{equation*}
    F_{\text{gate}} \approx 1 - \frac{1 \mu s}{100 \mu s} = 0.99
\end{equation*}
This means the gate has a fidelity of approximately 99\%. If we reduce the gate time to 0.5 $\mu s$, the fidelity improves to:
\begin{equation*}
    F_{\text{gate}} \approx 1 - \frac{0.5 \mu s}{100 \mu s} = 0.995
\end{equation*}
This means the gate has a fidelity of approximately 99.5\%. Therefore, optimizing gate times is crucial for enhancing the fidelity of quantum operations, as it directly reduces the error probability associated with decoherence.

\highspace
\textcolor{Green3}{\faIcon[regular]{lightbulb} \textbf{Threshold Theorem.}} Every physical quantum gate introduces a small probability of error.
Without error correction, these errors accumulate and eventually destroy the computation.

\highspace
The \definition{Threshold Theorem} states that if the physical \textbf{error rate} of the quantum hardware is \textbf{below a certain threshold value} (equivalently, if the gate \hl{fidelity is sufficiently high}), then \textbf{quantum error correction can remove errors faster than they are generated}.

\highspace
As a consequence, arbitrarily long quantum computations become possible using logical qubits and fault-tolerant protocols, even though the underlying physical qubits are noisy and have finite decoherence times.

\highspace
For example, if every second 10 errors appear, and we can correct 20 errors per second, then the system stays healthy. But if every second 10 errors appear, and we can only correct 5 errors per second, then the system becomes unhealthy and eventually fails. Therefore, maintaining an error rate below the threshold is crucial for the viability of quantum computing.